\chapter{系统定义与适用范围}

\section{系统定义}

Nous Wave 是面向长期运行人工智能主体的认知与人格平台。系统维护跨模型、会话、进程和检索索引持续存在的主体认知，使主体能够保留经历、修订认知、维护自我与社会关系，并在后续推理和行为中按来源、权限和上下文重新取得这些认知。

主体由稳定的逻辑身份及其认知状态定义。模型、模型服务商、会话和进程是可替换的执行条件；模型更换不得改变主体标识，也不得使持久认知失去来源谱系、修订历史或当前状态。

Nous Wave 负责认知语义与认知权威。通用数据库、索引、向量检索、对象存储、模型调用和协议传输由成熟基础设施承担；这些机制可以替换，只要不改变本设计规定的对象身份、领域归属、写入权限和可追溯性。

\section{能力组合}

同一套基础语义支持不同能力组合，不建立彼此分叉的产品实现。

\begin{ReportWideTable}{>{\raggedright\arraybackslash}p{3.3cm}>{\raggedright\arraybackslash}p{5.2cm}>{\raggedright\arraybackslash}p{5.5cm}}
\captionof{table}{Nous Wave 的主要能力组合}\label{tab:profiles}
\toprule
\textbf{能力组合} & \textbf{组成} & \textbf{用途}\\
\midrule
记忆型能力组合 & 主体基础、记忆、认知运行时、认知查询、检索服务、外部资源接口 & 为普通智能体、检索增强生成或其他调用方提供长期、可修订、可追溯的主体记忆。\\
完整主体能力组合 & 记忆型能力组合，加上自我、社会认知、动机和跨领域离线认知 & 维护长期主体的自我理解、社会世界、目标与离线认知演化。\\
\bottomrule
\end{ReportWideTable}

高级认知领域建立在同一主体基础、证据模型、修订机制和查询体系之上。未启用某个高级领域时，其缺失不得削弱记忆型能力组合的完整性。

\section{主体与模型}

主体标识（SubjectId）是跨限界上下文保持一致的逻辑身份。模型实例、模型名称、供应商账户、聊天线程、登录会话和宿主进程均不得作为主体身份。

模型输出只能以明确的角色进入系统：生成派生表示、形成候选提案、执行重排或生成表达。模型输出本身不具有认知权威。任何持久认知写入都由相应领域所有者按照当前对象状态、证据、权限和并发条件决定。

\section{设计抽象边界}

系统架构固定语义所有权、对象关系、生命周期边界、运行时边界、查询语义、持久化不变量以及与 Heptalogos 的系统边界。具体算法参数、数据库表、接口字段和部署细节可以在保持这些语义约束的前提下独立演化。
